
%%%%%%%%%%%%%%%%%%%%%%%%%%%%%%%%%%%%%%%%%%%%%%%%%%%%%%%%%%%%%
\subsubsection{具体分析}

品类级补货需同时解决需求预测和价格弹性两个子问题。需求预测方面，初始阶段采用ARIMA进行建模\cite{Wang2022}，经auto\_arima自动搜索，花叶类的最优阶数为(2,0,2)，辣椒类为(1,0,1)。然而ARIMA框架难以显式建模蔬菜销售的周周期结构与中国节假日效应——这些效应在生鲜零售场景中至关重要。Prophet模型通过显式分解趋势分量、周/年季节分量（基于傅里叶级数）以及中国节假日效应（含调休窗口），在此场景下具有明显的结构优势。价格弹性方面，文献中的弹性估计值存在较大分歧：杨涵等（2024）报告的弹性为$-1.92$，而陈妙霞等（2024）报告为$-0.46$\cite{Yang2024,Chen2024}，差异高达4倍。本文决定从自身数据估计各品类弹性，避免文献值与本数据集不匹配的风险。

首先对每个品类的日销售量序列进行Prophet分解，分离趋势、季节性和节假日效应，并从数据自行估计需求价格弹性。然后在需求预测和弹性估计的基础上，建立单周期报童模型，以加成率和订货量为联合决策变量，通过二维网格搜索求解最优策略。

\subsubsection{模型准备}

问题二以品类为单位，需将单品级数据按品类聚合。对每个品类$i$，计算每日总销售量$S_{i,t}$ = 销量加权平均售价$P_{i,t}$和加权平均批发价$C_{i,t}$。剔除批发价$\leq 0.01$元/kg的27条异常记录（占批发价数据的0.05\%）。各品类平均损耗率从附件4可见Sheet1的单品数据与附件1关联后按品类取均值计算：花菜类14.14\%，水生根茎类11.97\%，花叶类10.28\%，辣椒类8.52\%，食用菌8.13\%，茄类7.12\%。

需求价格弹性通过两阶段工具变量法从数据估计。Ordinary Least Squares估计弹性面临内生性问题：价格由商超在观测到需求信号后决定，价格与需求扰动相关，OLS估计量有偏且不一致。本文以批发价$C_{i,t}$作为工具变量——批发价由供应商决定，与当日的需求扰动外生，但通过成本加成定价与售价$P_{i,t}$强相关（第一阶段回归$R^2 > 0.7$）。两阶段估计为：
\begin{equation}
  \text{第一阶段：}\;\ln P_{i,t} = \beta_0 + \beta_1 \ln C_{i,t} + \sum_{k}\gamma_k D_{k,t} + \eta_{i,t}
\end{equation}
\begin{equation}
  \text{第二阶段：}\;\ln S_{i,t} = \beta_0' + e_i \widehat{\ln P}_{i,t} + \sum_{k}\gamma_k' D_{k,t} + \varepsilon_{i,t}
\end{equation}
其中$D_{k,t}$包括星期和月份虚拟变量，$\widehat{\ln P}_{i,t}$为第一阶段预测值。同时，对Prophet残差$\hat{\varepsilon}_{i,t}$建立局部加权回归（LOESS），捕捉价格响应中的非线性成分，避免Double-log的函数形式约束。估计结果如表\ref{tab:elasticity}所示，IV估计的弹性方向与OLS一致，但$R^2$从$0.011$--$0.144$提升至$0.25$--$0.43$，表明工具变量有效降低了遗漏变量偏误。

\begin{table}[H]
  \centering
  \caption{各品类需求价格弹性估计}
  \label{tab:elasticity}
  \begin{tabularx}{\textwidth}{CCCC}
    \toprule
    品类 & 价格弹性$e_i$（IV） & $R^2$（IV） & OLS $R^2$ \\
    \midrule
    水生根茎类 & $-1.01$ & 0.43 & 0.144 \\
    花叶类 & $-0.23$ & 0.25 & 0.019 \\
    花菜类 & $-0.73$ & 0.37 & 0.098 \\
    茄类 & $-0.40$ & 0.28 & 0.032 \\
    辣椒类 & $-0.38$ & 0.31 & 0.072 \\
    食用菌 & $-0.28$ & 0.25 & 0.011 \\
    \bottomrule
  \end{tabularx}
\end{table}

各品类弹性全部为负（$-0.23$至$-1.01$），符合需求定律。IV估计中$R^2$介于$0.25$--$0.43$，较OLS的$0.011$--$0.144$显著提升，验证了工具变量对遗漏变量偏误的修正效果。以花叶类为例，弹性$-0.23$意味着价格上升$10\%$仅导致销量下降约$2.3\%$，在实际运营中该效应可能被日常波动淹没。

将需求预测与价格弹性分阶段处理的合理性在于：Prophet已提取趋势、季节、节假日等主效应；工具变量法将批发价变动作为外生信号，隔离了价格内生性带来的偏误；残差中的价格响应通过LOESS非参数建模捕捉，不依赖Double-log的函数形式假设。灵敏度分析（表\ref{tab:sensitivity}）证实弹性扰动对最优利润影响$<0.1\%$，表明该策略的稳健性可接受。

该问题本质上是一个随机优化问题，需在需求不确定下联合优化加成率与订货量。最初考虑过将价格作为Prophet回归量以实现需求--价格联合估计，但价格与需求存在双向因果关系（内生性问题），导致OLS估计偏误。最终采用工具变量法分阶段处理：批发价作为工具变量识别弹性的外生成分，Prophet提取时间序列结构给出基准需求，LOESS残差建模捕获非线性的价格响应。灵敏度分析（表\ref{tab:sensitivity}）证实弹性扰动对最优利润的影响$<0.1\%$，表明分阶段策略的稳健性可接受。本文选择Prophet+约束报童作为主方法，ARIMA+固定加成作为baseline。
